\documentclass[10pt,a4paper]{article}
\usepackage[T1]{fontenc}
\usepackage[utf8]{inputenc}
\usepackage{lmodern}
\usepackage{amsmath,amssymb}
\usepackage[margin=21mm]{geometry}
\usepackage{microtype}
\usepackage[hidelinks]{hyperref}
\usepackage{enumitem}
\setlist[itemize]{leftmargin=*,nosep,topsep=3pt}
\setlength{\parindent}{0pt}
\setlength{\parskip}{5pt}
\setlength{\emergencystretch}{3em}
\newcommand{\question}[1]{\par\medskip{\normalsize\bfseries #1}\par\smallskip}
\pagestyle{plain}

\begin{document}
{\Large\bfseries The Development and Application of Optogenetics\par}

\textbf{Paper:} Fenno, Yizhar and Deisseroth (2011).\quad Review.\par
\href{https://doi.org/10.1146/annurev-neuro-061010-113817}{DOI: \nolinkurl{10.1146/annurev-neuro-061010-113817}}\par
{\footnotesize Page references below use PDF page numbers, including front matter.}

\question{To which general question does this paper relate?}
What optical stimulation can control, how targeting is achieved, and which actuator dynamics constrain a neural control model.

\question{How does it answer the general questions?}
\textbf{Actuation.} Reviews excitatory channelrhodopsins, inhibitory pumps and light-sensitive signalling tools. Genetic and projection targeting allow selective manipulation; opsin kinetics and spectral sensitivity vary substantially across tools (PDF pp.~2--11).

\textbf{Usable insight.} Light input is filtered by expression, delivery and opsin physiology before affecting spikes or circuit activity. High-frequency stimulation fidelity, adaptation and recovery therefore need validation for the selected construct. Electrical or optical measurements are needed to verify the effect; the actuator is not itself a sensor (PDF pp.~7--16).

\textbf{Model relevance.} Helps choose actuator location, sign and bandwidth and motivates adding opsin states when a static gain is inadequate. It does not establish the notebook's monotone memoryless nonlinearity assumptions, numerical delayed coupling coefficients or an LFP observation map. Discussed behavioural interventions are application-specific; this review does not identify beta suppression as a general therapeutic objective.

\question{What papers does it cite that could provide more information?}
\begin{itemize}
\item Boyden et al. (2005), \emph{Millisecond-timescale, genetically targeted optical control of neural activity}: optical excitation (references, PDF p.~19).
\item Gradinaru et al. (2009), \emph{Optical deconstruction of parkinsonian neural circuitry}: pathway-level intervention (PDF p.~20).
\end{itemize}

\end{document}
